\documentclass[12pt]{article}
\usepackage[utf8]{inputenc}
\usepackage{amsmath, amssymb, amsfonts}
\usepackage{graphicx}
\usepackage{hyperref}
\usepackage{enumerate}
\usepackage{float}
\usepackage{xcolor}
\usepackage{listings}
\usepackage{geometry}
\usepackage{fancyhdr}
\usepackage{mdframed}

\geometry{a4paper, margin=1in}

\newenvironment{enroutequestion}[1]{%
    \begin{mdframed}[linewidth=1pt, linecolor=blue!60, backgroundcolor=blue!5,
                     innertopmargin=8pt, innerbottommargin=8pt,
                     innerleftmargin=10pt, innerrightmargin=10pt]
    \noindent\textbf{Check-in Question #1.}~\itshape
}{%
    \end{mdframed}
}
\hypersetup{colorlinks=true, linkcolor=blue, urlcolor=blue}

\pagestyle{fancy}
\fancyhf{}
\lhead{Lab 7: Autonomy}
\rhead{MIT 1.104 2026}
\cfoot{\thepage}

\begin{document}

\begin{center}
    \linespread{1.2}
    \Large{\textbf{Lab 7: Autonomy}}\\
    \large{1.104 Team}\\
    \normalsize{Department of Civil and Environmental Engineering \\ Massachusetts Institute of Technology}\\
\end{center}

\section*{Introduction}
In this lab, you will learn about Proportional-Integral-Derivative (PID) control, one of the most widely used control techniques in industry.  PID controllers are found in a variety of applications ranging from simple temperature control to complex motion systems in robotics.  By the end of this lab, you will understand how each component of the PID controller affects system behavior and how to properly tune the parameters for optimal performance.  The lab consists of three parts: 
\begin{enumerate}
    \item PID control of a double integrator system 
    \item PID control of a rocket model (double integrator with gravity) 
    \item PID control of a nonlinear pendulum system 
\end{enumerate}

\section{Theoretical Background}

\subsection{PID Control Fundamentals}
A PID controller continuously calculates an error value $e(t)$ as the difference between a desired setpoint (SP) and a measured process variable (PV), and applies a correction based on proportional, integral, and derivative terms.  The mathematical form of a PID controller is: 

\begin{equation}
u(t) = K_p e(t) + K_i \int_{0}^{t} e(\tau) d\tau + K_d \frac{de(t)}{dt}
\end{equation}

Where:
$u(t)$ is the control signal, $e(t)$ is the error (setpoint - measured value), $K_p$ is the proportional gain, $K_i$ is the integral gain, $K_d$ is the derivative gain 

\subsection{The Role of Each Component}

\subsubsection{Proportional Term ($K_p$)}
The proportional term produces an output proportional to the current error value: 
\begin{equation}
P_{out} = K_p e(t)
\end{equation}

\textbf{Effect}: 
\begin{itemize}
    \item Increases system responsiveness 
    \item Reduces rise time (time to reach setpoint) 
    \item Too high: causes overshoot and oscillation 
    \item Too low: makes the system sluggish and unable to reach the setpoint 
\end{itemize}

\subsubsection{Integral Term ($K_i$)}
The integral term accumulates error over time and helps eliminate steady-state error: 
\begin{equation}
I_{out} = K_i \int_{0}^{t} e(\tau) d\tau
\end{equation}

\textbf{Effect}:
\begin{itemize}
    \item Eliminates steady-state error 
    \item Helps overcome system biases 
    \item Too high: causes overshoot and oscillation 
    \item Too low: steady-state error may persist 
\end{itemize}

\subsubsection{Derivative Term ($K_d$)}
The derivative term predicts system behavior by calculating the rate of change of error: 
\begin{equation}
D_{out} = K_d \frac{de(t)}{dt}
\end{equation}

\textbf{Effect}:
\begin{itemize}
    \item Reduces overshoot and settling time 
    \item Improves stability 
    \item Too high: amplifies noise and causes instability 
    \item Too low: system may exhibit excessive overshoot 
\end{itemize}

\subsection{Double Integrator System}
A double integrator system represents a basic model for many physical systems, like a mass moving in a frictionless environment.  Its mathematical representation is: 

\begin{equation}
\frac{d^2x}{dt^2} = u
\end{equation}

Where $x$ is the position, and $u$ is the control input (force).  In state-space form: 
\begin{equation}
\begin{bmatrix} \dot{x}_1 \\ \dot{x}_2 \end{bmatrix} = 
\begin{bmatrix} 0 & 1 \\ 0 & 0 \end{bmatrix}
\begin{bmatrix} x_1 \\ x_2 \end{bmatrix} + 
\begin{bmatrix} 0 \\ 1 \end{bmatrix} u
\end{equation}

Where $x_1$ is the position and $x_2$ is the velocity.  

\subsection{Rocket Model (Double Integrator with Gravity)}
The rocket model extends the double integrator by adding the effect of gravity, which creates a constant external force.  This model represents systems like a vertical rocket or an elevator, where gravity consistently affects the dynamics.  Its mathematical representation is: 

\begin{equation}
\frac{d^2x}{dt^2} = u - g
\end{equation}

Where $x$ is the position, $u$ is the control input (force), and $g$ is the gravitational acceleration.  In state-space form: 
\begin{equation}
\begin{bmatrix} \dot{x}_1 \\ \dot{x}_2 \end{bmatrix} = 
\begin{bmatrix} 0 & 1 \\ 0 & 0 \end{bmatrix}
\begin{bmatrix} x_1 \\ x_2 \end{bmatrix} + 
\begin{bmatrix} 0 \\ 1 \end{bmatrix} u + 
\begin{bmatrix} 0 \\ -g \end{bmatrix}
\end{equation}

The key difference from the simple double integrator is that the constant gravitational force creates a bias that must be compensated for by the controller.  This makes the integral term ($K_i$) particularly important, as it accumulates error over time and provides the necessary constant control input to counteract gravity.  

\subsection{Pendulum System}
The pendulum is described by the nonlinear differential equation: 

\begin{equation}
ml^2\ddot{\theta} = mgl\sin\theta - b\dot{\theta} + u
\end{equation}

Where:
$\theta$ is the angle from the upright position (0 is upright, $\pi$ is downward), $m$ is the mass at the end of the pendulum, $l$ is the length of the pendulum, $g$ is the gravitational acceleration, $b$ is the damping coefficient, $u$ is the control torque 

This system is nonlinear due to the $\sin\theta$ term, making control more challenging.  Note that in this formulation, the gravity term has a positive sign, as it acts to increase the angle away from the upright position.  

\subsection{Linear Quadratic Regulator (LQR)}

The Linear Quadratic Regulator (LQR) is an optimal control technique that provides a systematic way to design feedback controllers for linear systems.  Unlike PID control, which requires manual tuning, LQR determines control gains mathematically to minimize a specified performance index.  

Each experiment includes a ground truth solution using LQR.  This provides a benchmark for comparison with your PID controller. LQR is an advanced control method that provides optimal performance.  Use it as a reference when tuning your PID controller. 

\section{Performance Metrics for Control Systems}

Before beginning the experiments, it's important to understand the key metrics used to evaluate control system performance.  Throughout all experiments in this lab, you'll need to record and analyze these metrics: 

\begin{itemize}
    \item \textbf{Rise Time}: The time it takes for the system output to rise from 10\% to 90\% of the steady-state value.  A shorter rise time indicates a more responsive system but may lead to overshoot.  
    \item \textbf{Overshoot}: The amount by which the system output exceeds the setpoint, expressed as a percentage of the setpoint value.  For example, if the setpoint is 1.0 and the maximum value reached is 1.2, the overshoot is 20\%.  
    \item \textbf{Settling Time}: The time required for the system response to reach and stay within a specified range (typically $\pm$2\% or $\pm$5\%) of the final value.  This indicates how quickly the system stabilizes. 
    \item \textbf{Steady-State Error}: The difference between the desired setpoint and the final value of the system output after transients have died out.  Ideally, this should be zero or very close to zero.  
\end{itemize}

Each script \textbf{automatically computes and displays} these four metrics directly on the response plot, along with visual markers (rise-time point, peak overshoot, settling-time line, and a $\pm$2\% settling band).  A summary text box in the corner of the plot shows the numerical values, so you can read and record them without manual calculation.

You should record these metrics under different controller configurations and analyze how changes to the PID parameters affect them.  These observations will form the core of your lab report.  Note that it's not necessary to record these metrics for \textbf{each run}, but you should record the metrics for several runs to get a sense of the performance.  

\section{Laboratory Procedure}

\subsection{Setup}
Before beginning the lab, ensure you have the necessary software installed.  You need \textbf{Python 3.8 or newer} (check with \texttt{python --version}).

\begin{enumerate}
    \item On the lab laptop, open \texttt{VS Code}, and press \texttt{Ctrl+`} to open the terminal. 
    \item Run the following command to move to the \texttt{Desktop} folder:
    \begin{verbatim}
    cd Desktop
    \end{verbatim}
    \item Clone the lab repository:
    \begin{verbatim}
    git clone https://github.com/LiJiarui111/MIT-1104-Lab7-PID-2025
    \end{verbatim}
    \item Navigate into the repository and install the required Python packages from the provided \texttt{requirements.txt}:
    \begin{verbatim}
    cd MIT-1104-Lab7-PID-2025
    pip install -r requirements.txt
    \end{verbatim}
    This installs \texttt{numpy}, \texttt{scipy}, and \texttt{matplotlib}.
    \item Verify the installation by running one of the scripts:
    \begin{verbatim}
    python double_integrator.py
    \end{verbatim}
    Two windows should appear: a line plot and an animation.  Close them to return to the terminal.
    \item Read the \texttt{README.md} file for additional setup details and troubleshooting tips.
\end{enumerate}

\textbf{Note on running scripts.}  Throughout this lab you will run three scripts.  The commands are:
\begin{verbatim}
    python double_integrator.py
    python rocket_model.py
    python pendulum.py
\end{verbatim}
If \texttt{python} is not recognised on your system, try \texttt{python3} instead.

\subsection{General Tuning Protocol}
For \textbf{each} of the three systems below, you are required to test and record the performance of \textbf{5 different parameter sets}. Open the corresponding Python script in VS Code and change the \texttt{Kp}, \texttt{Ki}, and \texttt{Kd} variables at the top of the file (inside the clearly marked block), save the file (\texttt{Ctrl+S}), then run the script in the terminal.  The performance metrics will be displayed automatically on the plot---record them in your lab report table.

The 5 sets must follow this progression:
\begin{enumerate}
    \item \textbf{Set 1 (P-Only):} Set $K_i = 0$ and $K_d = 0$. Tune $K_p$ to achieve a reasonable rise time.
    \item \textbf{Set 2 (Add I or D):} Keep your $K_p$ from Set 1. Add \textit{either} $K_i > 0$ or $K_d > 0$ to observe how it alters steady-state error or transient overshoot. 
    \item \textbf{Set 3 (Full PID):} Use non-zero values for $K_p$, $K_i$, and $K_d$.
    \item \textbf{Set 4 (Free Tuning):} Experiment freely to improve your metrics relative to Set 3.
    \item \textbf{Set 5 (Best Performance):} Your final, highly-tuned parameter set attempting to match or beat the optimal LQR baseline shown in the plots.
\end{enumerate}

\subsection{Part 1: Double Integrator System}
Open \texttt{double\_integrator.py} in VS Code.  In the terminal, run:
\begin{verbatim}
    python double_integrator.py
\end{verbatim}
Two windows will appear: a \textbf{response plot} (with annotated performance metrics) and an \textbf{animation}.  Close both windows to return to the terminal, then change the PID gains and re-run.  Follow the general tuning protocol to complete your 5 sets. 

Once you have a good PID parameter set for the double integrator, you can go to the TA and get them tested on an Arduino robot car.  It's very interesting to see how the PID controller works on a real autonomous system.  

\begin{enroutequestion}{1}
Run the double integrator with P-only control ($K_i = 0$, $K_d = 0$) using a moderate $K_p$.  Now double $K_p$.  What happens to the rise time and the overshoot?  Based on your observations, explain in one or two sentences why increasing $K_p$ alone cannot simultaneously minimise both rise time and overshoot.
\end{enroutequestion}

\subsection{Part 2: Rocket Model System}
Open \texttt{rocket\_model.py} in VS Code.  In the terminal, run:
\begin{verbatim}
    python rocket_model.py
\end{verbatim}
Follow the general tuning protocol. Pay special attention to how the ground truth (LQR) solution handles the gravity disturbance.  Note how the control signal must overcome gravity to maintain position.  The annotated plot will show you how large the steady-state error is when $K_i = 0$---this is a key observation for the rocket system.

\begin{enroutequestion}{2}
Run the rocket model with $K_i = 0$ and observe the steady-state error.  Now set $K_i$ to a small positive value (e.g.\ 1.0) and re-run.  Why does the integral term reduce the steady-state error for this system but was less critical for the double integrator?  (Hint: think about the role of the constant gravitational force $g$.)
\end{enroutequestion}

\begin{enroutequestion}{3}
Compare your best Set~5 results for the double integrator and the rocket model.  Which system was harder to tune, and why?  Did you find that the same $K_p$ and $K_d$ values that worked well for the double integrator also worked well for the rocket, or did you need to change them significantly?
\end{enroutequestion}

\subsection{EXTRA: Part 3: Pendulum System}
Open \texttt{pendulum.py} in VS Code.  In the terminal, run:
\begin{verbatim}
    python pendulum.py
\end{verbatim}
Follow the general tuning protocol. Observe the system response and note whether the pendulum successfully reaches the upright position.  Watch the animation that shows both the PID-controlled pendulum and the ground truth (LQR) solution.  

\begin{enroutequestion}{4}
The pendulum equation contains a $\sin\theta$ nonlinearity, whereas the double integrator and rocket are linear systems.  Try using the exact same PID gains that worked well for the double integrator.  Does the pendulum still behave well?  In one or two sentences, explain why tuning a PID controller for a nonlinear system is generally more difficult than for a linear one.
\end{enroutequestion}

\section{Additional Report Requirements}

In your lab report, please include some discussion about PID parameters tuning, such as:
\begin{itemize}
    \item A table showing your 5 parameter sets and their corresponding rise time, overshoot, settling time, and steady-state error for each system.  (These values are displayed on the plot after each run---you can read them directly from the summary text box.)
    \item Your answers to the four \textbf{Check-in Questions} embedded in the procedure above.
    \item How do you determine the best parameters?  
    \item What is the effect of each parameter on the system response?  
    \item Comparison between double integrator and pendulum (if completed). 
    \item Challenges encountered and solutions. 
\end{itemize}

\section{References}

\begin{enumerate}
    \item Åström, K. J., \& Hägglund, T. (2006). Advanced PID control.  ISA-The Instrumentation, Systems, and Automation Society.
    \item Ogata, K. (2010). Modern control engineering (5th ed.). Prentice Hall.  
    \item Dorf, R. C., \& Bishop, R. H. (2011). Modern control systems (12th ed.). Pearson. 
\end{enumerate}

\end{document}